% RESEARCH_DRAFT_SHA256: 44c228d8bb46ed0b4af406b03c27ca35281ca44aa45b2af7c31b461a1b5cb77b
% Layout and source-reference wrappers expanded for the web renderer.
\documentclass[11pt]{article}
\usepackage[margin=25mm]{geometry}
\usepackage{fontspec,unicode-math}
\setmainfont{Latin Modern Roman}
\setmathfont{Latin Modern Math}
\usepackage{amsmath,amssymb,mathtools,xurl,hyperref}
\hypersetup{colorlinks=true,urlcolor=blue,linkcolor=blue}
\setcounter{secnumdepth}{0}
\setlength{\parindent}{0pt}
\setlength{\parskip}{5pt}
\setlength{\emergencystretch}{3em}
\newcommand{\sref}[2]{\hyperlink{source:#1:#2}{[S#2]}}
\newcommand{\cataloguescope}[1]{\paragraph{Source record.}#1}
\newcommand{\floor}[1]{\left\lfloor#1\right\rfloor}
\begin{document}
\section*{Research-attempt provenance}
This standalone research attempt retains the complete original definition body below,
including its historical source and preparation notes. Those notes describe the
earlier catalogue preparation; the new attempt and its verification are identified
separately after that body. No original ID, catalogue status, or source statement is changed.

\section{What is actually exponential in single-carrier equalization?}\label{statement}

\subsection{Definitions and mathematical statement}

\paragraph{Shared notation.}Standard number systems are $\mathbb N,\mathbb Z,\mathbb Q,\mathbb R,\mathbb C$. Specialized notation and index ranges follow the cited source and the definitions below.


For a discrete channel $y_t=\sum_{j=0}^{N-1}h_jx_{t-j}+z_t$, equalization seeks to undo intersymbol interference. An OFDM scheme uses orthogonal subcarriers and FFT-based transforms; a cyclic prefix can convert the relevant convolution into a circular one under stated conditions. Maximum-likelihood sequence detection can have a state count exponential in channel memory, whereas linear equalization solves a different optimization problem.
\paragraph{Record scope and qualifications.}
The catalogue title is retained as metadata. Named certificates, contracts, bounds or frameworks in that title are not asserted as results of the original source extract.
The quoted exponential complexity is not a claim about every single-carrier equalizer. Algorithm, alphabet size, channel memory, and accuracy must be distinguished. The composite appendix text is preserved literally. The supplied catalogue marks this record partially solved; its present research status was not independently checked.
The source text below is rendered literally where necessary; damaged formulas and composite extracts are not silently repaired into a new mathematical claim.
\paragraph{Statement recorded in the supplied source.}
1. Equalization complexity AP: MC modulation, e.g., OFDM, can handle the equalization task easily. However, for HDB channels, one would need long symbols, and problems with channel variability and channel estimation arise. Equalization complexity in MC modulation is O(N log N ), whereas in SC modulation, the complexity is exponential in N (N

is the channel length). AP: Linear vs. nonlinear equalization Linear equalization algorithms are not efficient. It appears that nonlinear techniques, e.g., THP, DPC, perform better, but their interactions with SF have not been investigated. AP: Special equalization techniques The problem frequently reduces to that of equalizing a long but sparse channel. These channels are encountered in TV systems, and are handled in DVB-T with OFDM-like techniques. 2. OFDM AP: OFDM is the system of choice for 3.5G, 4G and Wi-Fi systems. The delay spread and the Doppler define the cyclic prefix, symbol length and FFT size for an OFDM system. BH: Coding in OFDM with fine frequency spacing, such as in HDB channels that are very frequency selective, becomes more complex. AP: Spatial focusing OFDM and SF have not been studied. We expect some perfor-mance and complexity trade-offs. 3. Waveform Selection TS: The waveform selection problem is not only relevant to HDB chan-nels. AP: In UWB systems, OFDM and PPM are contenders. 4. Open vs. closed loop techniques 16 AP: What happens when you don't know the channel exactly? What happens when you estimate the channel on the uplink and use this knowledge on the downlink? 17 A Acronyms Acronyms

CDMA Code Division Multiple Access CSI Channel State Information DFE Decision Feedback Equalizer DPC Dirty Paper Coding DVB-T Digital Video Broadcasting - Terrestrial FDD Frequency Division Duplexing FFT Fast Fourier Transform FIR Finite Impulse Response HDB High Delay spread Bandwidth HSDPA High Speed Data Packet Access IIR Infinite Impulse Response ISI Inter Symbol Interference ITU International Telecommunication Union LE Linear Equalizer LMS Least Mean Square LPI Low Probability of Intercept MC Multi Carrier MIMO Multiple Input Multiple Output MISO Multiple Input Single Output ML Maximum Likelihood 18 OFDM Orthogonal Frequency Division Multiplexing PDP Power Delay Profile PPM Pulse Position Modulation QAM Quadrature Amplitude Modulation Rx Receiver SC Single Carrier SDMA Space Division Multiple Access SF Spatial Focusing SIMO Single Input Multiple Output SINR Signal to Interference plus Noise Ratio SNR Signal to Noise Ratio SONAR Sonic Radar TDD Time Division Duplexing THP Tomlinson Harashima Precoding TR Time Reversal TV Television Tx Transmitter UWB Ultra Wideband 19 B Participants Abbreviations of Names of Participants

AA Agissilaos Athanassoulis AF Albert Fannjiang AK Arnold Kim AP Arogyaswami Paulraj BH Babak Hassibi CT Chrysoula Tsogka DB Deborah Berebichez FL Frederick Lee ES Erik Stauffer GD Gregoire Derveaux GL Gabriel Lerosey GP George Papanicolaou HB Harry Bims HS Heechun Song JF Jean-Pierre Fouque JH Jan Hansen JS James Frederick Sifferlen KS Knut Solna LB Liliana Borcea MF Mathias Fink 20 MV Mai Vu MC Mohamad Charafeddine ME Majid Emami OO Ozgur Oyman PK Persefoni Kyritsi RC Robert Calderbank RD Robert Clark Daniels RH Robert Heath RM Raghuraman Mudumbai RN Rohit Nabar ST Stephan ten Brink TS Thomas Strohmer UM Upamanyu Madhow 21 \href{https://aimath.org/WWN/timerev/timerev.pdf}{[S2]}
\paragraph{Source status.}
The supplied catalogue labels this source record partially solved; that metadata was checked against the supplied snapshot on 2026-10-11. The label is preserved for provenance and is not an independent verification of present open status or mathematical correctness.
\subsection{Short English statement}
The source compares equalization methods and waveform choices, then includes acronym and participant lists.
\subsection{Sources}
\begin{itemize}
\item[S1] Ulam AI and the UnsolvedMath Contributors, UnsolvedMath: A Curated Collection of Open Mathematics Problems, catalogue entry. Catalogue curation and metadata: CC BY 4.0. The source extract is retained; definitions, English summaries and qualifications are editorial additions.
\url{https://huggingface.co/datasets/ulamai/UnsolvedMath}.
\item[S2] AIM Workshop Problem Lists, source linked by the catalogue.
\url{https://aimath.org/WWN/timerev/timerev.pdf}.
\end{itemize}
\paragraph{Preparation and review.}
The original source and UnsolvedMath contributors retain attribution for the source extract. Codex (GPT-6; exact serving revision not exposed) assisted with definitions, English summaries, scope notes and formatting. No new human mathematical authorship or independently verified solution is asserted. Catalogue identity checked against the complete supplied registry and current repository. Distinct upstream records, including repeated questions and broad or damaged extracts, are retained for complete coverage of the source CS category. A repeated record is not a new independent conjecture. Complete mathematical-equivalence and current-literature review has not been performed. Preparation and final review follow contribution prompts C and D, with source gaps disclosed. Structural and rendering checks do not establish a new theorem or open status.
\label{end:statement}

\section{Literature and context check}
The equalization discussion and acronym list in the primary AIM time-reversal
PDF were consulted. It compares several different receiver objectives without
fully specifying one complexity problem. The calculation below fixes exact
finite-alphabet least-squares sequence estimation to test the exponential claim.
\section{Research attempt: A long sparse channel need not require exponential-state equalization}
Let $x_t$ lie in a known alphabet of size $q$, with known preceding symbols,
and observations
$y_t=\sum_{j=0}^m h_jx_{t-j}+z_t$ for $1\le t\le T$.
For independent equal-variance Gaussian noise, maximizing likelihood is
equivalent to minimizing the sum of squared residuals. A dynamic program with
state the last $m$ symbols has $q^m$ states; for each new symbol it adds the
current residual and shifts the state. The minimum cost for a state is sufficient
because all future residuals depend on the past only through those $m$ symbols.
There are at most $q^{m+1}$ transitions per time, yielding
$O(Tq^{m+1}(m+1))$ arithmetic work if each residual is evaluated directly.
Precomputed or incrementally evaluated transition residuals can alter polynomial
factors but this displayed generic implementation has exponential state count.

The attempted inference was that channel length alone forces such complexity.
Take a sparse channel with only $h_0$ and $h_d$ nonzero, for an arbitrarily large
delay $d$. Its objective is
\[
 \sum_{t=1}^T|y_t-h_0x_t-h_dx_{t-d}|^2.
\]
Terms couple symbols only in the same residue class modulo $d$. Partition all
indices by that residue. The objective is a sum of independent chain objectives,
each with a nearest-neighbor interaction and known initial boundary symbol.
Each chain is minimized by a $q$-state dynamic program, testing $q^2$
successive-symbol pairs at each position. Across all chains there are $T$
positions, so the exact global minimizer costs $O(Tq^2)$ arithmetic operations,
plus initialization of the nonempty chains. Backtracking reconstructs the
complete symbol sequence. This is exact discrete sequence estimation, not a
linear relaxation.

The sparse support therefore matters through the interaction structure, not
only the largest delay. Adding a tap at delay one couples different residue
chains and invalidates the decomposition; no general sparse-channel polynomial
bound follows. Conversely the generic exponential-state implementation is not
a lower bound for all algorithms. Channel-estimation uncertainty and other
coding/precoding objectives in the source remain distinct problems.
\section{Verification and adversarial audit}
The chain factorization was verified term by term modulo $d$, including
$d>T$ (each observed symbol then has a known predecessor) and $d=1$ (one
ordinary chain). Initial symbols were specified, and no cross-chain coding
constraint is allowed in the restricted calculation.
\paragraph{Audit corrections.} An exponential dynamic-program state count was not asserted to be an unconditional lower bound, and sparse delay structure was tested explicitly.
\paragraph{External computation.} None; symbolic arguments and boundary checks only.
\section{Scope and remaining gaps}
The first gap is the exact input, output, channel-support and coding model
for the intended complexity comparison. The remaining OFDM, focusing,
waveform and unknown-channel questions are not solved by the two-tap case.
\paragraph{Integration review.} Codex (GPT-6; exact serving revision not exposed) reviewed the written arguments and expanded layout and source-reference wrappers. Repository metadata and the production web renderer were checked using Python 3.13 and JavaScript with MathJax 3.2.2. This model-assisted review is not independent human proof verification; no PDF compilation or formal proof checking is asserted.

\section{Completion Estimate}
\noindent\textbf{COMPLETION: 0\%}
The calculations below the preserved source statement investigate a concrete route,
but establish no substantiated new progress toward the entire original target.
Known special cases and failed sufficient conditions are not credited as a solution.
\section{References}
The original source references and catalogue attribution are preserved in the
Sources section above. The following URLs were actually consulted for this attempt;
retention of an additional original reference does not claim fresh consultation of it.
\begin{itemize}
\item \url{https://aimath.org/WWN/timerev/timerev.pdf} (consulted 2026-10-11). 
\end{itemize}
\section{Model and contributor disclosure}
OpenAI gpt-6-astra, requested reasoning configuration ultra, generated this research
attempt on 2026-10-11. The exact serving revision is not exposed. Codex GPT-6
performed orchestration, repository-format review, and a model-assisted mathematical
audit of the written arguments; its exact serving revision
is also not exposed. The mathematical route, boundary checks and final adversarial
self-audit followed contribution prompts A then B in
\texttt{docs/contribution\_prompts.md}, after reading \texttt{CONTRIBUTING.md}.
Python 3.14.2 assembled the standalone documents; any mathematical computation is
specified in Verification. No theorem prover was used. No human mathematical
authorship, substantial human proof correction, or independent human proof checking
is asserted. Model-assisted review, source retrieval, and structural checks do not
establish correctness or novelty. Full private reasoning transcripts are unavailable
as reproducibility artifacts; the complete arguments, inputs and outcomes actually
used are stated here or in the companion batch audit.
\section{Final status}
Unresolved as a new research contribution. This is a concrete route investigation,
not a claim of new complete resolution or a current-status certification. A known
answer reproduced for context does not become a new contribution. The catalogue's
source-status label is preserved; this attempt status does not assert that a
source-labelled solved question is open.
\end{document}
